\documentclass[11pt]{amsart}

\usepackage[T1]{fontenc}
\usepackage{lmodern}
\usepackage{microtype}
\usepackage{mathtools,amssymb,amsthm}
\usepackage{array}
\usepackage[hidelinks]{hyperref}

\hypersetup{
  pdftitle={Polynomial Certificates for the First Two Partial Quotients of
            sqrt((zeta(3)-S_N)/2), and the Third for N at least 2}
}

\newtheorem{theorem}{Theorem}
\newtheorem{lemma}[theorem]{Lemma}
\newtheorem{proposition}[theorem]{Proposition}
\newtheorem{corollary}[theorem]{Corollary}
\theoremstyle{definition}
\newtheorem*{conjecture}{Conjecture (Feldman)}
\theoremstyle{remark}
\newtheorem*{remark}{Remark}

\DeclareMathOperator{\PQ}{a}
\newcommand{\Rt}{R}
\newcommand{\ps}{\psi}
\newcommand{\pss}{\psi_3}

\title[Certificates for leading partial quotients]
{Polynomial Certificates for the First Two Partial Quotients of
 $\sqrt{(\zeta(3)-S_N)/2}$, and the Third for $N\ge2$}

\date{}

\subjclass[2020]{11A55, 11J70, 11Y65, 11M06}
\keywords{continued fraction, partial quotient, Riemann zeta at three,
quasi-polynomial, Euler--Maclaurin, positivity certificate, Pell number}

\begin{document}

\begin{abstract}
Let $S_N=\sum_{j=1}^{N}j^{-3}$, let $\Rt_N=\zeta(3)-S_N$ and let
$g(N)=\sqrt{\Rt_N/2}$. Feldman displays the continued fraction
$g(N)=\bigl[0;2N+1,\,4N+2,\,\lfloor(12N+6)/19\rfloor,\ldots\bigr]$ and records, as
the first of his open problems, that his uniform Euler--Maclaurin remainder does
not pin even the first two of these at the smallest index; in particular his
claim that $\PQ_1(g)|_{N_k}=Q_{2k+1}$ and $\PQ_2(g)|_{N_k}=2Q_{2k+1}$ along the
Pell indices $N_k$ is left unproved. We prove $\PQ_1(g(N))=2N+1$ and
$\PQ_2(g(N))=4N+2$ for every integer $N\ge1$, which gives that claim at every
$k\ge1$. We also prove $\PQ_3(g(N))=\lfloor(12N+6)/19\rfloor$ for every $N\ge2$;
at $N=1$ the displayed value $0$ is not the third partial quotient of $g(1)$,
which is $1$. Each floor decision is reduced to the non-negativity of the
coefficient vector of an explicit integer polynomial; no value of $\zeta(3)$ is
used in any proof. The trailing ellipsis of the display is untouched, so
Feldman's conjecture as a whole remains open.
\end{abstract}

\maketitle

\section{The statement}

For a real $x$ we write $x=[a_0;\PQ_1,\PQ_2,\ldots]$ for its regular continued
fraction and $\PQ_i(x)$ for its $i$-th partial quotient. Throughout,
\[
 S_N=\sum_{j=1}^{N}\frac1{j^3},\qquad
 \Rt_N=\zeta(3)-S_N=\sum_{j>N}\frac1{j^3},\qquad
 g(N)=\sqrt{\frac{\Rt_N}{2}} .
\]
Since $\Rt_1=\zeta(3)-1<2$ we have $0<g(N)<1$, so $a_0=0$ and the index $1$
really is the first floor-and-invert. Also $g(N)$ is irrational for every $N$:
$\zeta(3)$ is irrational \cite{Apery,Beukers} and $S_N$ is rational, so
$\Rt_N\notin\mathbb Q$, and a rational $g(N)$ would make
$\Rt_N=2g(N)^2$ rational. Hence the continued fraction of $g(N)$ is infinite
and every partial quotient below is unambiguous.

Feldman \cite{Feldman} introduces $g$ and states the following, quoted from the
e-print source.

\begin{conjecture}[\cite{Feldman}]
The formal continued fraction of $g(N)=\sqrt{\Rt_N/2}$ is
\begin{equation}\label{eq:gcf}
 g(N)=\Bigl[0;\;2N+1,\;4N+2,\;
   \Bigl\lfloor\frac{12N+6}{19}\Bigr\rfloor,\;\ldots\Bigr]
\end{equation}
with quasi-polynomial partial quotients of periods $1,1,19,\ldots$. In
particular, at $N=N_k$:
\begin{equation}\label{eq:gcfpell}
 \PQ_1(g)\big|_{N_k}=2N_k+1=Q_{2k+1},\qquad
 \PQ_2(g)\big|_{N_k}=4N_k+2=2Q_{2k+1}.
\end{equation}
\end{conjecture}

Here $P_k$ are the Pell numbers $0,1,2,5,12,29,70,169,\dots$, $Q_k=P_k+P_{k-1}$,
$M_k=P_{2k+1}$ and $N_k=(Q_{2k+1}-1)/2=3,20,119,696,\ldots$, so that
$Q_{2k+1}^2-2M_k^2=-1$. Thus \eqref{eq:gcfpell} is the restriction of the first
two entries of \eqref{eq:gcf} to $N=N_k$, and is strictly weaker than
\eqref{eq:gcf}.

What is open is stated by the author himself, at the first of his open problems,
again quoted:
\begin{quote}
``The derivation is only formal: the uniform Euler--Maclaurin remainder
\dots\ does not pin the leading floors at the smallest index
$k=1$ ($M_k=5$), so even the first two quotients
$\PQ_1(g)|_{N_k}=Q_{2k+1}$, $\PQ_2(g)|_{N_k}=2Q_{2k+1}$ are not proved there.
A sharper index-dependent remainder, or a direct treatment of the small cases,
would settle it.''
\end{quote}
The author's main theorem, which is proved, concerns $\sqrt{\Rt_{N_k}}$ --- a
different real number, $\sqrt2$ times ours --- and only for $k\ge3$, so it says
nothing at $k=1$.

We prove the following.

\begin{theorem}\label{thm:a12}
For every integer $N\ge1$,
\[
 \PQ_1\bigl(g(N)\bigr)=2N+1
 \qquad\text{and}\qquad
 \PQ_2\bigl(g(N)\bigr)=4N+2 .
\]
\end{theorem}

\begin{theorem}\label{thm:a3}
For every integer $N\ge2$,
\[
 \PQ_3\bigl(g(N)\bigr)=\Bigl\lfloor\frac{12N+6}{19}\Bigr\rfloor ,
\]
and this fails at $N=1$: there $\PQ_3(g(1))=1$ while
$\lfloor 18/19\rfloor=0$.
\end{theorem}

Theorem~\ref{thm:a12} discharges \eqref{eq:gcfpell} at every $k\ge1$, the
index $k=1$ ($N_1=3$, $M_1=5$) included, which is exactly the statement the
author's open problem names as unproved. Theorem~\ref{thm:a3} is an additional
numeric statement about the third entry of \eqref{eq:gcf}, with $N=1$ excluded.
What is \emph{not} settled is recorded in Section~\ref{sec:status}.

The mechanism is elementary, which is why no value of $\zeta(3)$ appears: we
build two short antidifferences of opposite defect sign, obtaining a tail-free
two-sided enclosure of $\Rt_N$ valid from $N=1$; each floor decision then becomes
the positivity of a rational function of $N$ whose shifted numerator has
non-negative integer coefficients.

\section{Two antidifferences of opposite defect sign}\label{sec:anti}

Put
\begin{equation}\label{eq:psi}
 \ps(x)=\frac1{2x^2}+\frac1{2x^3}+\frac1{4x^4}-\frac1{12x^6}
        =\frac{A(x)}{12x^{6}},
 \qquad
 \pss(x)=\ps(x)+\frac1{12x^8}=\frac{B(x)}{12x^{8}},
\end{equation}
where
\begin{equation}\label{eq:psi3}
 A(x)=6x^4+6x^3+3x^2-1,
 \qquad
 B(x)=x^2A(x)+1=6x^6+6x^5+3x^4-x^2+1.
\end{equation}
Both vanish at infinity. Define the defects
\[
 D(j)=\ps(j)-\ps(j+1)-\frac1{j^3},
 \qquad
 D_3(j)=\pss(j)-\pss(j+1)-\frac1{j^3}.
\]

\begin{lemma}\label{lem:signs}
As identities of polynomials in $j$,
\begin{align}
 12\,j^6(j+1)^6\,D(j)   &= -(2j+1)^3=-\bigl(8j^3+12j^2+6j+1\bigr),
 \label{eq:D}\\
 12\,j^8(j+1)^8\,D_3(j) &= 18j^5+45j^4+48j^3+27j^2+8j+1 .
 \label{eq:D3}
\end{align}
Hence $D(j)<0$ and $D_3(j)>0$ for every $j\ge1$.
\end{lemma}

\begin{proof}
Both are polynomial identities after clearing denominators, and both are
verified by expansion. Explicitly, \eqref{eq:D} is
$A(j)(j+1)^6-A(j+1)j^6-12j^3(j+1)^6=-(2j+1)^3$, and \eqref{eq:D3} is
$B(j)(j+1)^8-B(j+1)j^8-12j^5(j+1)^8=18j^5+45j^4+48j^3+27j^2+8j+1$. The two
coefficient vectors $(-1,-6,-12,-8)$ and $(1,8,27,48,45,18)$, in ascending
powers of $j$, have all entries of one sign, and the signs are opposite; the
denominators $12j^6(j+1)^6$ and $12j^8(j+1)^8$ are positive for $j\ge1$. As a
hand check at $j=1$: $\ps(1)-\ps(2)-1=\tfrac76-\tfrac{155}{768}-1=-\tfrac{27}{768}$
and $\pss(1)-\pss(2)-1=\tfrac{49}{1024}$.
\end{proof}

\begin{lemma}\label{lem:enclosure}
For every integer $N\ge1$,
\begin{equation}\label{eq:enc}
 \ps(N+1)\;<\;\Rt_N\;<\;\pss(N+1),
\end{equation}
and moreover
\begin{equation}\label{eq:enctight}
 \pss(N+1)-\frac{49}{40N^{10}}\;<\;\Rt_N\;<\;\pss(N+1).
\end{equation}
\end{lemma}

\begin{proof}
Since $\ps(x)\to0$, the series $\sum_{j>N}\bigl(\ps(j)-\ps(j+1)\bigr)$
telescopes exactly to $\ps(N+1)$, whence
\[
 \Rt_N=\sum_{j>N}\frac1{j^3}
      =\sum_{j>N}\Bigl(\ps(j)-\ps(j+1)\Bigr)-\sum_{j>N}D(j)
      =\ps(N+1)-\sum_{j>N}D(j),
\]
and likewise $\Rt_N=\pss(N+1)-\sum_{j>N}D_3(j)$. Both tail series converge
absolutely, since \eqref{eq:D} and \eqref{eq:D3} give
$|D(j)|=O(j^{-9})$ and $D_3(j)=O(j^{-11})$. By Lemma~\ref{lem:signs} the
first tail is negative and the second positive, which is \eqref{eq:enc}; note
that the upper bound carries \emph{no} remainder term.

For \eqref{eq:enctight} we bound the second tail. From \eqref{eq:D3},
$D_3(j)\le \tfrac{49}{4}j^{-11}$ for every $j\ge1$: clearing denominators, the
claim is that $588(j+1)^8-4j^3\bigl(18j^5+45j^4+48j^3+27j^2+8j+1\bigr)$ has
non-negative coefficients, and indeed it equals
$516j^8+4524j^7+16272j^6+32820j^5+41128j^4+32924j^3+16464j^2+4704j+588$. Next,
$\sum_{j>N}j^{-11}\le\tfrac1{10}N^{-10}$, because
$\tfrac1{10}j^{-10}-\tfrac1{10}(j+1)^{-10}\ge(j+1)^{-11}$ for $j\ge1$ --- which
after clearing denominators reads
$(j+1)^{11}-j^{11}-11j^{10}\ge0$, true coefficientwise --- and the left side
telescopes. Hence $0<\sum_{j>N}D_3(j)\le\tfrac{49}{4}\cdot\tfrac1{10}N^{-10}$.
\end{proof}

\section{The first two quotients}\label{sec:a12}

Write $V=1/g(N)=\sqrt{2/\Rt_N}$ and
\begin{equation}\label{eq:lohi}
\begin{aligned}
 \ell(N)&=(2N+1)+\frac1{4N+3}=\frac{8N^2+10N+4}{4N+3},\\
 h(N)&=(2N+1)+\frac1{4N+2}=\frac{8N^2+8N+3}{4N+2}.
\end{aligned}
\end{equation}

\begin{proof}[Proof of Theorem~\ref{thm:a12}]
Consider the two rational functions
\[
 I_2(N)=h(N)^2\,\ps(N+1)-2,
 \qquad
 I_1(N)=2-\ell(N)^2\,\pss(N+1).
\]
Clearing denominators, $I_2=\mathcal I_2/\bigl(12(N+1)^6(4N+2)^2\bigr)$ and
$I_1=\mathcal I_1/\bigl(12(N+1)^8(4N+3)^2\bigr)$ with
$\mathcal I_1,\mathcal I_2\in\mathbb Z[N]$; substituting $N=1+t$,
\begin{equation}\label{eq:I2}
 \mathcal I_2(1+t)=659+1380\,t+1055\,t^2+342\,t^3+38\,t^4,
\end{equation}
\begin{equation}\label{eq:I1}
\begin{aligned}
 \mathcal I_1(1+t)={}&492+3448\,t+7664\,t^2+8328\,t^3\\
 &{}+5032\,t^4+1728\,t^5+316\,t^6+24\,t^7 .
\end{aligned}
\end{equation}
Every coefficient in \eqref{eq:I2} and \eqref{eq:I1} is positive, and both
denominators are manifestly positive for $N\ge1$; therefore
$I_1(N)>0$ and $I_2(N)>0$ for every integer $N\ge1$. (As a hand check
of \eqref{eq:I1} at $N=1$: $2\cdot12\cdot256\cdot49-484\cdot621=301056-300564=492$.
Values of $\mathcal I_2$ at $N=1,\dots,5$ are $659,3474,10983,26606,54675$; of
$\mathcal I_1$, $492,27032,263772,1415016,5415116$.)

Now $I_2(N)>0$ says $h(N)^2>2/\ps(N+1)$, and $\Rt_N>\ps(N+1)$ by
\eqref{eq:enc}, so $h(N)^2>2/\ps(N+1)>2/\Rt_N=V^2$ and hence $V<h(N)$.
Symmetrically $I_1(N)>0$ says $\ell(N)^2<2/\pss(N+1)$, and
$\Rt_N<\pss(N+1)$, so $\ell(N)^2<2/\pss(N+1)<2/\Rt_N=V^2$ and $V>\ell(N)$.
Thus
\begin{equation}\label{eq:Vbox}
 (2N+1)+\frac1{4N+3}\;<\;V\;<\;(2N+1)+\frac1{4N+2}
 \qquad(N\ge1).
\end{equation}
Since $0<\tfrac1{4N+2}<1$, \eqref{eq:Vbox} gives $\lfloor V\rfloor=2N+1$, i.e.
$\PQ_1(g(N))=2N+1$. Taking reciprocals of the fractional part in
\eqref{eq:Vbox}, $1/(V-(2N+1))$ lies strictly between $4N+2$ and $4N+3$, so its
floor is $4N+2$, i.e. $\PQ_2(g(N))=4N+2$.
\end{proof}

\begin{corollary}\label{cor:pell}
$\PQ_1(g)|_{N_k}=Q_{2k+1}$ and $\PQ_2(g)|_{N_k}=2Q_{2k+1}$ for every $k\ge1$;
that is, \eqref{eq:gcfpell} holds, including at $k=1$.
\end{corollary}

\begin{proof}
$N_k=(Q_{2k+1}-1)/2$ gives $2N_k+1=Q_{2k+1}$ and $4N_k+2=2Q_{2k+1}$
identically; apply Theorem~\ref{thm:a12} at $N=N_k\ge3$.
\end{proof}

\section{The third quotient}\label{sec:a3}

The displayed third quotient is quasi-polynomial of period $19$: writing
$N=19s+\rho$ with $0\le\rho\le18$,
\begin{equation}\label{eq:A3}
 A_3(N):=\Bigl\lfloor\frac{12N+6}{19}\Bigr\rfloor
 =12s+\Bigl\lfloor\frac{12\rho+6}{19}\Bigr\rfloor .
\end{equation}
Fix $\rho$ and put, with $A_1=2N+1$ and $A_2=4N+2$ as above,
\[
 \ell_3=A_1+\cfrac1{A_2+\cfrac1{A_3}},
 \qquad
 h_3=A_1+\cfrac1{A_2+\cfrac1{A_3+1}} ,
\]
which are rational functions of $s$ with $\ell_3<h_3$. Because
$\PQ_1=A_1$ and $\PQ_2=A_2$ are already known from Theorem~\ref{thm:a12},
$\PQ_3(g(N))=A_3$ holds if and only if $\ell_3<V<h_3$. (Both $\ell_3$ and $h_3$
require $A_3\ge1$; every class below is certified only from an $s_0$ at which
$A_3\ge1$, the extreme case being $\rho=0$, $s_0=1$, $A_3=12$.) Exactly as in
Section~\ref{sec:a12}, using \eqref{eq:enctight} in place of \eqref{eq:enc},
$\ell_3<V$ follows from the positivity of
\[
 J_1(s)=2-\ell_3^{\,2}\,\pss(N+1)\Big|_{N=19s+\rho},
\]
and $V<h_3$ from the positivity of
\[
 J_2(s)=h_3^{\,2}\Bigl(\pss(N+1)-\frac{49}{40N^{10}}\Bigr)-2\Big|_{N=19s+\rho}.
\]
(No separate check that the lower bound in \eqref{eq:enctight} is positive is
needed: if it were not, $J_2$ would be at most $-2$.)

\begin{proposition}\label{prop:classes}
For each of the $19$ residue classes $\rho$, the numerators and denominators of
both $J_1$ and $J_2$, written as polynomials in $s-s_0$, have all coefficients
non-negative, with $s_0=s_0(\rho)$ as tabulated:
\begin{center}
\normalfont
\setlength{\tabcolsep}{2.5pt}
\footnotesize
\begin{tabular}{l|ccccccccccccccccccc}
 $\rho$ & 0&1&2&3&4&5&6&7&8&9&10&11&12&13&14&15&16&17&18\\
 \hline
 $\lfloor(12\rho+6)/19\rfloor$
         & 0&0&1&2&2&3&4&4&5&6&6&7&7&8&9&9&10&11&11\\
 $s_0(\rho)$
         & 1&1&1&1&1&0&0&0&0&0&0&0&0&0&0&0&0&0&0\\
\end{tabular}
\end{center}
Consequently $\PQ_3(g(N))=A_3(N)$ for every $N\ge5$, and the only integers
$N\ge1$ not covered are $N\in\{1,2,3,4\}$.
\end{proposition}

\begin{proof}
Each entry is a finite exact computation with integer-coefficient polynomials in
one variable: substitute $N=19s+\rho$ and $A_3=12s+\lfloor(12\rho+6)/19\rfloor$,
clear denominators, shift $s\mapsto s_0+t$ and read the signs. As the covering
statement: the classes $\rho=5,\dots,18$ certify from $s=0$, giving
$N=5,\dots,18$ and then $N=19s+\rho$ for all $s\ge1$; the classes
$\rho=0,1,2,3,4$ certify from $s=1$, giving $N=19,20,21,22,23$ and upwards.
The union is $\{N: N\ge5\}$, and its complement in $\mathbb Z_{\ge1}$ is
$\{1,2,3,4\}$.
\end{proof}

In the class $\rho=9$ one has $12N+6\equiv0\pmod{19}$, so the floor in
\eqref{eq:A3} sits \emph{on} an integer; that class too certifies from $s=0$,
with
\begin{multline}\label{eq:rho9}
 \operatorname{num}J_1\big|_{\rho=9}
 =12491451+142067432\,s+671815989\,s^2+1690280508\,s^3\\
 {}+2385550453\,s^4+1789910922\,s^5+557513238\,s^6 ,
\end{multline}
of degree $6$ with all coefficients positive, while $\operatorname{num}J_2$ has
$18$ positive coefficients, the constant one being $37139362880053056000$. Here
$\operatorname{num}$ is the numerator over
$12(N+1)^{8}\bigl(A_2A_3+1\bigr)^{2}$, evaluated at $N=19s+9$ with
$A_2=4N+2$ and $A_3=12s+6$; as a check on the normalisation, the right-hand side
of \eqref{eq:rho9} takes the values $12491451$, $7249629993$ and
$147632938319$ at $s=0,1,2$.

\section{The four remaining indices}\label{sec:small}

The four indices $N\in\{1,2,3,4\}$ left over by
Proposition~\ref{prop:classes} are decided by hand from
Lemma~\ref{lem:enclosure} together with an exact rational partial sum
$\sum_{j=N+1}^{J}j^{-3}$, for one small $J$. Since
$\Rt_N=\sum_{j=N+1}^{J}j^{-3}+\Rt_J$, applying \eqref{eq:enc} at $J$ gives
\begin{equation}\label{eq:small}
\begin{aligned}
 \tfrac{155}{768}&<\Rt_1<\tfrac{207}{1024}
   &&(J=1), &\quad g(1)&=[0;3,6,1,\ldots],\\
 \tfrac{337}{4374}&<\Rt_2<\tfrac{6067}{78732}
   &&(J=2), &\quad g(2)&=[0;5,10,1,\ldots],\\
 \tfrac{1967}{49152}&<\Rt_3<\tfrac{10491}{262144}
   &&(J=3), &\quad g(3)&=[0;7,14,2,\ldots],\\
 \tfrac{1707247}{69984000}&<\Rt_4<\tfrac{61461017}{2519424000}
   &&(J=5), &\quad g(4)&=[0;9,18,2,\ldots].
\end{aligned}
\end{equation}
In each line the printed bracket already determines the first three partial
quotients: one applies $x\mapsto1/x$ and $x\mapsto x-\lfloor x\rfloor$ three
times to the two endpoints and observes that the two integer parts agree at
every step. This completes the proof of Theorem~\ref{thm:a3}: the displayed
values at $N=2,3,4$ are $1,2,2$, which \eqref{eq:small} confirms, while at
$N=1$ the displayed value is $\lfloor18/19\rfloor=0$ and the true value is $1$.

\begin{remark}[the exceptional index]
$\lfloor(12N+6)/19\rfloor=0$ can only happen at $N=1$, and $0$ is not a legal
partial quotient at a position $\ge1$; so under the numeric reading the third
entry of \eqref{eq:gcf} needs the hypothesis $N\ge2$. That reading may not be
the intended one: \eqref{eq:gcf} is captioned ``the \emph{formal} continued
fraction'', and under a strictly formal, series-level reading it makes no
assertion at the integer $N=1$.
\end{remark}

\section{Scope}\label{sec:status}

\noindent\textbf{Not settled.} The trailing ellipsis of
\eqref{eq:gcf}. Nothing here addresses $\PQ_4$, $\PQ_5$ or their
quasi-periods, and the phrase ``periods $1,1,19,\ldots$'' therefore remains a
conjecture: \textbf{the Conjecture of Section~1 is not proved}. Nor is anything
proved about $\PQ_i(g(N))$ for $i\ge4$, about the growth of the quotients, or
about irrationality or normality questions for $\zeta(3)$.

\medskip\noindent\textbf{No Lean.} We wrote no Lean and inspected none: what is
discharged is a mathematical statement, not a machine proof.

\medskip\noindent\textbf{Attribution.} The values $2N+1$, $4N+2$ and
$\lfloor(12N+6)/19\rfloor$ are Feldman's, printed in \cite{Feldman}, which also
observes that \eqref{eq:gcfpell} is immediate from $N_k=(Q_{2k+1}-1)/2$ once the
quotients are known.

\medskip\noindent\textbf{Reproduction.} Everything above is finite and exact.
Lemmas~\ref{lem:signs} and~\ref{lem:enclosure}, Theorem~\ref{thm:a12} and
Proposition~\ref{prop:classes} are polynomial identities and coefficient-sign
readings in $\mathbb Z[N]$ or $\mathbb Z[s]$; Section~\ref{sec:small} is four
rational brackets. No floating-point number, no enumeration, no solver and no
decimal expansion of $\zeta(3)$ enters any proof. The accompanying program
\texttt{verify.py} re-derives every displayed quantity from the definitions
\eqref{eq:psi}--\eqref{eq:psi3} alone, in exact integer and \texttt{Fraction}
arithmetic and using only the Python standard library, and reports $95$ passing
checks; the recorded run \texttt{verify.output.txt} also states, in the
program's own words, what it does not cover.

\begin{thebibliography}{9}

\bibitem{Feldman}
D.~V. Feldman,
\emph{Anomalous partial quotients in the continued fraction of
$\sqrt{\zeta(3)-S_N}$},
arXiv:2607.04077v1, 2026.
\href{https://arxiv.org/abs/2607.04077}{arXiv:2607.04077}.
The conjecture and open problem quoted in Section~1 are the statements labelled
\texttt{prop:gcf} and \texttt{op:gcf} in the e-print source \texttt{q2.tex};
the main theorem referred to there is \texttt{thm:main}, its uniform remainder
\texttt{lem:key}, and the corresponding Lean declaration \texttt{prop\_gcf}.
The e-print was at v1 and had no journal version when this was written.

\bibitem{Apery}
R.~Ap\'ery,
\emph{Irrationalit\'e de $\zeta(2)$ et $\zeta(3)$},
Ast\'erisque \textbf{61} (1979), 11--13.

\bibitem{Beukers}
F.~Beukers,
\emph{A note on the irrationality of $\zeta(2)$ and $\zeta(3)$},
Bull. London Math. Soc. \textbf{11} (1979), no.~3, 268--272.
\href{https://doi.org/10.1112/blms/11.3.268}{doi:10.1112/blms/11.3.268}.

\end{thebibliography}

\end{document}
